\documentclass[11pt]{amsart}
\usepackage[T1]{fontenc}
\usepackage{lmodern}
\usepackage[margin=1in]{geometry}
\usepackage{amsmath,amssymb,amsthm,mathtools,microtype,etoolbox}
\usepackage[colorlinks=true,linkcolor=blue,urlcolor=blue,citecolor=blue]{hyperref}
\hypersetup{
  pdftitle={A sharp Hardy inequality on Z2},
  pdfsubject={Sharp logarithmic Hardy inequality and positive remainder},
  pdfauthor={Nika Areshidze}
}
\newtheorem{theorem}{Theorem}[section]
\newtheorem{lemma}[theorem]{Lemma}
\newtheorem{proposition}[theorem]{Proposition}
\newtheorem{corollary}[theorem]{Corollary}
\numberwithin{equation}{section}
\newcommand{\Z}{\mathbb Z}
\newcommand{\R}{\mathbb R}
\newcommand{\E}{\mathcal E}
\newcommand{\V}{\mathcal V}
\newcommand{\supp}{\operatorname{supp}}
\newcommand{\Cc}{C_c}
\allowdisplaybreaks[1]
\title[A sharp Hardy inequality on $\Z^2$]{A sharp Hardy inequality on $\Z^2$}
\author[N. Areshidze]{Nika Areshidze}
\address{Department of Mathematics, University of California, Irvine, CA 92697, USA}
\email{nareshid@uci.edu}
\keywords{Discrete Hardy inequality, square lattice, sharp constant,
supersolution, logarithmic weight}
\date{November 4, 2026}

% Display the date below the author name.
\makeatletter
\patchcmd{\@adminfootnotes}
  {\ifx\@empty\@date\else \@footnotetext{\@setdate}\fi}
  {}{}{\errmessage{Unable to relocate the date}}
\patchcmd{\@setauthors}
  {\MakeUppercase{\authors}}
  {\MakeUppercase{\authors}\par\nobreak\vspace{6pt}%
   {\normalfont\normalsize\@date\par}}
  {}{\errmessage{Unable to display the date below the author}}
\makeatother

\begin{document}
\begin{abstract}
We prove the Hardy inequality on $\Z^2$ with Euclidean weight
$|x|^{-2}(\log|x|)^{-2}$ and sharp constant $1/4$, for finitely supported
functions vanishing on $\{|x|\le1\}$.
\end{abstract}
\maketitle
\vspace{-18pt}

\section{Introduction}

On the exterior domain $\Omega=\{x\in\R^2:|x|>1\}$, the classical
logarithmic Hardy inequality
\[
\int_{\Omega}|\nabla u(x)|^2\,dx
\ge\frac14\int_{\Omega}
\frac{|u(x)|^2}{|x|^2(\log|x|)^2}\,dx
\]
holds for every $u\in C_c^\infty(\Omega)$. The constant $1/4$ is sharp.
We establish the corresponding inequality for the nearest-neighbor
Laplacian on $\Z^2$, with the same Euclidean weight and the same sharp
constant, for finitely supported functions vanishing on $\{|x|\le1\}$.

Throughout, $|x|$ denotes the Euclidean norm and $\log$ the natural
logarithm. Let $e_1=(1,0)$, $e_2=(0,1)$, and set
\[
K=\{x\in\Z^2:|x|\le1\}=\{0,\pm e_1,\pm e_2\},
\qquad D=\Z^2\setminus K.
\]
We identify $\Cc(D)$ with the real-valued finitely supported functions
on $\Z^2$ that vanish on $K$. Their Dirichlet energy is
\begin{equation}\label{eq:energy}
\E(u)=\sum_{x\in\Z^2}\sum_{j=1}^2|u(x+e_j)-u(x)|^2,
\end{equation}
so that each unoriented edge is counted once.

\begin{theorem}\label{thm:main}
Every $u\in\Cc(D)$ satisfies
\begin{equation}\label{eq:main}
\sum_{x\in\Z^2}\sum_{j=1}^2|u(x+e_j)-u(x)|^2
\ge\frac14\sum_{x\in D}\frac{|u(x)|^2}{|x|^2(\log|x|)^2}.
\end{equation}
The constant $1/4$ is sharp, and the inequality is strict for every
nonzero $u\in\Cc(D)$.
\end{theorem}

Kapitanski and Laptev \cite[Theorem 4.1]{KL} established a logarithmic
Hardy inequality on $\Z^2$ with an unspecified positive coefficient.
Their theorem uses the same Dirichlet set $K$, with the Euclidean weighted
sum restricted to $|x|\ge2$. Inequality \eqref{eq:main} also includes
the four points $(\pm1,\pm1)$.

Huang and Ye \cite[Remark 6.2]{HY} obtained a logarithmic Hardy weight
on $\Z^2$ with leading coefficient $1/4$, for functions supported on
$\{|x|\ge2\}$. Its next term is negative along diagonal directions,
so their expansion does not establish \eqref{eq:main} with the
prescribed weight.

Hake, Keller and Pogorzelski \cite[Corollary 2 and the subsequent
remark]{HKP} constructed an optimal Hardy weight from the lattice
potential kernel. Its leading term is
$1/[4|x|^2(\log|x|)^2]$, followed by a negative correction of order
$|x|^{-2}(\log|x|)^{-3}$. They proved sharpness of the leading coefficient
near infinity and discussed whether $1/4$ remains valid globally for the
prescribed logarithmic weight. Theorem~\ref{thm:main} establishes this
endpoint for $K$. Their criticality result concerns the different class
of functions vanishing only at the origin.

The proof uses the ground-state identity
\cite{KPP}. We start from the radial function
$\sqrt{\frac12\log(|x|^2-\frac12)}$.
An angular comparison bounds its four-neighbor sum by the corresponding
sum in a diagonal direction, where an exact factorization gives the
required estimate. Changing the values on the two orbits $|x|^2=2,4$
then produces a global supersolution. The construction yields a positive
remainder, while logarithmic cutoffs prove sharpness.

A direct consequence is the following family of inequalities.

\begin{corollary}\label{cor:shift}
For each $b\ge1$, every $u\in\Cc(D)$ satisfies the sharp inequality
\begin{equation}\label{eq:shift}
\E(u)\ge\frac14\sum_{x\in D}
\frac{|u(x)|^2}{|x|^2\log^2(b|x|)}.
\end{equation}
\end{corollary}

The inequality follows from \eqref{eq:main} because
$\log(b|x|)\ge\log|x|$ on $D$; sharpness is proved in
Section~\ref{sec:sharpness}. The inequalities and their sharp constants
also hold for complex-valued functions, by applying the inequalities
to the real and imaginary parts.

\section{The supersolution}\label{sec:super}

Write $x\sim y$ when $|x-y|=1$, and define
\[
L\phi(x)=4\phi(x)-\sum_{y\sim x}\phi(y),
\qquad
W(x)=\frac1{4|x|^2(\log|x|)^2}\quad(x\in D).
\]
Thus $\E(u)=\sum_xu(x)Lu(x)$. Our goal is to construct $\phi>0$ on
$D$, with $\phi=0$ on $K$, such that $L\phi>W\phi$ on $D$.
The following identity explains this reduction.

\begin{lemma}\label{lem:ground}
If $\phi>0$ on $D$ and $\phi=0$ on $K$, then for every $u\in\Cc(D)$,
\begin{equation}\label{eq:ground}
\E(u)=
\sum_{\substack{\{x,y\}:x\sim y\\x,y\in D}}
\phi(x)\phi(y)
\left(\frac{u(x)}{\phi(x)}-\frac{u(y)}{\phi(y)}\right)^2
+\sum_{x\in D}\frac{L\phi(x)}{\phi(x)}u(x)^2.
\end{equation}
Each edge in the first sum is counted once.
\end{lemma}

\begin{proof}
For $x,y\in D$, expansion gives
\begin{align*}
(u(x)-u(y))^2
&=\phi(x)\phi(y)
\left(\frac{u(x)}{\phi(x)}-\frac{u(y)}{\phi(y)}\right)^2\\
&\quad+\frac{\phi(x)-\phi(y)}{\phi(x)}u(x)^2
+\frac{\phi(y)-\phi(x)}{\phi(y)}u(y)^2.
\end{align*}
Sum over the edges in $D$. An edge from $x\in D$ to $K$ contributes
$u(x)^2$ to the energy and the same amount to the potential term,
since $u=\phi=0$ on $K$. This proves \eqref{eq:ground}.
\end{proof}

\begin{lemma}\label{lem:angular}
Let $q>0$, and suppose that $F\in C^4$ on an open interval containing
$[q+1-2\sqrt q,q+1+2\sqrt q]$, with $F^{(4)}\le0$ there.
If $x_1^2+x_2^2=q$, then
\begin{align}
&\sum_{j=1}^2\bigl(F(q+1+2x_j)+F(q+1-2x_j)\bigr)\notag\\
&\hspace{15mm}\le
2F(q+1+\sqrt{2q})+2F(q+1-\sqrt{2q}).\label{eq:angular}
\end{align}
\end{lemma}

\begin{proof}
Put $a=q+1$ and
$G(s)=F(a+2\sqrt s)+F(a-2\sqrt s)$ for $0\le s\le q$.
For $s>0$, with $t=2\sqrt s$, differentiation gives
\[
G''(s)=\frac4{t^3}
\left[
t\bigl(F''(a+t)+F''(a-t)\bigr)
-\int_{a-t}^{a+t}F''(v)\,dv
\right]\le0.
\]
Indeed, $F''$ is concave and therefore lies above the chord joining
its endpoint values. By continuity, $G$ is concave on $[0,q]$.
Jensen's inequality now gives
$G(x_1^2)+G(x_2^2)\le2G(q/2)$, which is \eqref{eq:angular}.
\end{proof}

We use the function
\[
F(v)=\sqrt{\frac12\log\left(v-\frac12\right)}\qquad(v>3/2)
\]
and define
\begin{equation}\label{eq:phi}
\phi(x)=
\begin{cases}
0,&|x|^2\le1,\\[1mm]
\dfrac35,&|x|^2=2,\\[2mm]
\dfrac{18}{25},&|x|^2=4,\\[2mm]
F(|x|^2),&|x|^2\ge5.
\end{cases}
\end{equation}
These cases exhaust $\Z^2$, and $\phi>0$ on $D$.

\begin{proposition}\label{prop:super}
The function \eqref{eq:phi} satisfies $(L-W)\phi>0$ on $D$.
More precisely, for $q=|x|^2\ge8$,
\begin{equation}\label{eq:globalradial}
\frac{L\phi(x)}{\phi(x)}\ge\frac1{q\log^2(q-\frac12)}>W(x),
\end{equation}
and for $1<|x|^2<8$, we have $(L-W)\phi(x)>1/1250$.
\end{proposition}

\begin{proof}
Suppose first that $q\ge8$. Then $\phi(x)=F(q)$ and
$\phi(y)\le F(|y|^2)$ for every neighbor $y$ of $x$.
Indeed, a neighbor can meet a modified orbit only when $q=9$ and
$|y|^2=4$, and $18/25<F(4)$, as verified below.
Writing $t=\log(v-\frac12)>0$, we have
\[
F^{(4)}(v)=
-\frac{48t^3+44t^2+36t+15}
{16\sqrt2\,(v-\frac12)^4t^{7/2}}<0.
\]
Since $(\sqrt q-1)^2\ge(\sqrt8-1)^2>3/2$,
Lemma~\ref{lem:angular} applies to the four neighboring values.
Set
\[
h=\log\left(q-\frac12\right),\qquad
d=\log\frac{\sqrt{2q}+1}{\sqrt{2q}-1}.
\]
The factorization
$q+\frac12\pm\sqrt{2q}=\frac12(\sqrt{2q}\pm1)^2$
gives
\[
\log\left(q+\frac12\pm\sqrt{2q}\right)=h\pm d,
\qquad 0<d<h.
\]
The elementary inequalities
\[
\sqrt{1+z}+\sqrt{1-z}\le2-\frac{z^2}{4},
\qquad
\log\frac{1+z}{1-z}\ge2z
\qquad(0\le z<1)
\]
follow, respectively, by squaring and by integrating
$2/(1-z^2)\ge2$. In particular, $d\ge\sqrt{2/q}$.
The angular comparison therefore gives
\begin{align*}
\frac{L\phi(x)}{\phi(x)}
&\ge4-2\sqrt{1+d/h}-2\sqrt{1-d/h}\\
&\ge\frac{d^2}{2h^2}
\ge\frac1{qh^2}>\frac1{q(\log q)^2}=W(x).
\end{align*}
This proves \eqref{eq:globalradial}.

It remains to check $q=2,4,5$. Write $W_q=1/[q(\log q)^2]$.
We use the rational bounds
\begin{gather}
\frac{18}{25}<F(4),\qquad
\frac{173}{200}<F(5)<\frac{87}{100},\notag\\
F(8)<\frac{201}{200},\qquad F(9)<\frac{26}{25},\qquad
F(10)<\frac{213}{200},\label{eq:boundsF}\\
W_2<\frac{21}{20},\qquad W_4<\frac2{15},\qquad
W_5<\frac2{25}.\label{eq:boundsW}
\end{gather}
All these bounds follow by substitution in the elementary enclosure
\begin{equation}\label{eq:logcertificate}
0<\log t-2\sum_{k=0}^{4}\frac{z^{2k+1}}{2k+1}
<\frac{2z^{11}}{11(1-z^2)},
\qquad z=\frac{t-1}{t+1},\quad t>1.
\end{equation}
To obtain it, expand $\log t=2\sum_{k\ge0}z^{2k+1}/(2k+1)$
and bound the tail by a geometric series.
Substituting \eqref{eq:boundsF}--\eqref{eq:boundsW} into the
three neighbor sums gives
\begin{align}
(L-W)\phi(1,1)
&=(4-W_2)\frac35-2F(5)>\frac3{100},\notag\\
(L-W)\phi(2,0)
&=(4-W_4)\frac{18}{25}-F(9)-2F(5)>\frac1{250},\label{eq:local}\\
(L-W)\phi(2,1)
&=(4-W_5)F(5)-\frac35-\frac{18}{25}-F(8)-F(10)
>\frac1{1250}.\notag
\end{align}
Sign changes and coordinate permutations cover all remaining points.
\end{proof}

\begin{proof}[Proof of the lower bound in Theorem~\ref{thm:main}]
Apply \eqref{eq:ground} to the function \eqref{eq:phi}.
The first term is nonnegative, and Proposition~\ref{prop:super}
gives $L\phi/\phi>W$ at every point of $D$. Hence
$\E(u)\ge\sum_D W(x)u(x)^2$, with strict inequality for $u\ne0$.
\end{proof}

The same estimates give an explicit remainder.

\begin{corollary}\label{cor:remainder}
For $x\in D$, put $q=|x|^2$ and
\[
R(x)=
\begin{cases}
\dfrac1{1250},&1<q<8,\\[2mm]
\displaystyle\frac1{q\log^2(q-\frac12)}
-\frac1{q(\log q)^2},&q\ge8.
\end{cases}
\]
Then $R(x)>0$ on $D$, and every $u\in\Cc(D)$ satisfies
\begin{equation}\label{eq:remainder}
\E(u)\ge
\frac14\sum_{x\in D}\frac{|u(x)|^2}{|x|^2(\log|x|)^2}
+\sum_{x\in D}R(x)|u(x)|^2.
\end{equation}
\end{corollary}

\begin{proof}
For $1<q<8$, we have $\phi(x)<1$, so \eqref{eq:local} gives
$(L-W)\phi(x)/\phi(x)>1/1250$. For $q\ge8$, use
\eqref{eq:globalradial}. The conclusion follows from \eqref{eq:ground}.
\end{proof}

The coefficient $1/4$ is therefore sharp within the family of constant
multiples of the prescribed weight, although the weight admits a positive
additive improvement. No optimality is asserted for the remainder $R$.

\section{Sharpness}\label{sec:sharpness}

Choose $\chi\in C^\infty([0,\infty))$ with $0\le\chi\le1$,
$\chi=1$ on $[0,1]$, and $\chi=0$ on $[2,\infty)$.
For $T\ge2$, put
\[
u_T(x)=
\begin{cases}
\sqrt{\log|x|}\,\chi\bigl(T^{-1}\log|x|\bigr),&x\in D,\\
0,&x\in K.
\end{cases}
\]
These functions are finitely supported. Write
\[
\V_b(u)=\sum_{x\in D}
\frac{|u(x)|^2}{|x|^2\log^2(b|x|)}
\quad(b\ge1),\qquad \V=\V_1.
\]

Let $U_T(z)=\sqrt{\log|z|}\,\chi(T^{-1}\log|z|)$ for $|z|>1$.
On $r=|z|\ge2$, differentiation gives, uniformly in $T$,
\[
|\nabla U_T(z)|\le Cr^{-1},\qquad
|\nabla^2U_T(z)|\le Cr^{-2}.
\]
The gradients of $|\nabla U_T|^2$ and $U_T^2/(r^2\log^2r)$
are consequently bounded by $Cr^{-3}$.
Taylor's theorem and comparison on unit squares therefore give
uniformly bounded errors when replacing lattice sums by integrals
outside a fixed ball, since $\sum_{|x|\ge3}|x|^{-3}<\infty$.
The finitely many remaining lattice terms are also uniformly bounded.

With $t=\log r$, it follows that
\begin{align*}
\V(u_T)
&=2\pi\int_{\log2}^{\infty}
\frac{\chi(t/T)^2}{t}\,dt+O(1)
=2\pi\log T+O(1),\\
\E(u_T)
&=2\pi\int_{\log2}^{\infty}
\left(
\frac{\chi(t/T)}{2\sqrt t}
+\frac{\sqrt t}{T}\chi'(t/T)
\right)^2dt+O(1)\\
&=\frac\pi2\log T+O(1).
\end{align*}
Indeed, $\chi'=0$ unless $T\le t\le2T$, so the terms containing
$\chi'$ contribute $O(1)$. Hence
$\E(u_T)/\V(u_T)\to1/4$, proving sharpness in
Theorem~\ref{thm:main}.

For fixed $b\ge1$, the bound $u_T(x)^2\le\log|x|$ gives
\[
0\le\V(u_T)-\V_b(u_T)
\le2(\log b)\sum_{x\in D}
\frac1{|x|^2(\log|x|)^2}=O_b(1).
\]
The series converges, so $\E(u_T)/\V_b(u_T)\to1/4$,
proving Corollary~\ref{cor:shift}.
Removing any fixed finite set from the weighted sums changes them
by $O(1)$; the coefficient remains sharp after such a removal.

\begin{thebibliography}{9}

\bibitem{HKP}
P. Hake, M. Keller and F. Pogorzelski,
\emph{Optimal Hardy inequalities for random walks on $\Z^2$},
arXiv:2608.24213v1, 25 August 2026.
\url{https://arxiv.org/abs/2608.24213v1}.

\bibitem{HY}
X. Huang and D. Ye,
\emph{One-dimensional sharp discrete Hardy--Rellich inequalities},
J. Lond. Math. Soc. (2) \textbf{109} (2024), no.~1, e12838.
\url{https://doi.org/10.1112/jlms.12838}.

\bibitem{KL}
L. Kapitanski and A. Laptev,
\emph{On continuous and discrete Hardy inequalities},
J. Spectr. Theory \textbf{6} (2016), no.~4, 837--858.
\url{https://doi.org/10.4171/JST/144}.

\bibitem{KPP}
M. Keller, Y. Pinchover and F. Pogorzelski,
\emph{Optimal Hardy inequalities for Schr\"odinger operators on graphs},
Comm. Math. Phys. \textbf{358} (2018), no.~2, 767--790.
\url{https://doi.org/10.1007/s00220-018-3107-y}.

\end{thebibliography}
\end{document}
